% Supplementary figures of the article, adapted by a one-off script from the earlier draft
% legends in figure_captions_draft.tex (panel letters as (a), cross-references to the article
% through xr as M-<label>). Each figure is a macro, placed by supplementary_erl.tex in its
% section. The article's audit figure is erl_fig_audit (the ERL version of fig4_audit, without its
% per-block maps, which are the two block figures below).
% Edit this file, not the draft, for the Supplementary.

\newcommand{\SIfiglevels}{%
\begin{figure}[!htbp]
  \centering
  \includegraphics[width=\linewidth]{../figures/out/fig1_levels.pdf}
  \caption{\textbf{Basin water level, scene record and epochs.}
  (a) Daily basin level on the DAHITI reference (normal heights, not the Baltic Height System), interpolated between the \metric{levels.dahiti.n_records} altimetry records (points)\citep{schwatke2015}, with Hydroweb levels shifted onto that reference where records are more than \metric{params.max_gap_days}~days apart (orange)\citep{cretaux2011,hydrowebwl}. Shading and dashed lines, the epochs E1 (to \metric{levels.epoch_E1_end}), E2 (to \metric{params.e2_end}) and E3. The median 1$\sigma$ of a scene level, \metric{levels.scene_sigma_p50_m}~m, is not drawn.
  (b) and (c) Levelled scenes per calendar month: Landsat~5, 7, 8 and 9 (b)\citep{usgs_l45tm_c2l2,usgs_l7etm_c2l2,usgs_l89oli_c2l2} and Sentinel-2 granules (c; modified Copernicus Sentinel data), $n = \metric{levels.n_scenes_by_epoch.E1}$ (E1), \metric{levels.n_scenes_by_epoch.E2} (E2) and \metric{levels.n_scenes_by_epoch.E3} (E3) scenes before screening. Blue, open-water months; grey, the excluded ice season (\metric{params.ice_first_month} to \metric{params.ice_last_month}).}
  \label{fig:levels}
\end{figure}
}

\newcommand{\SIfigstudyarea}{%
\begin{figure}[!htbp]
  \centering
  \includegraphics[width=\linewidth]{../figures/out/fig_study_area.pdf}
  \caption{\textbf{Study domain, reporting sectors, computation cover and validation blocks.}
  (a) The study domain (thick line; \metric{study_area.domain_km2}~\kmsq{}, land included) and its seven reporting sectors (thin lines, numbered as in the key), frozen before any evaluation. Grey, the cover over which water frequencies are computed (\metric{study_area.stage2_cover_km2}~\kmsq{}); light blue, open sea at the \metric{params.lowstand_label} low stand.
  (b) The $n = \metric{blocks.n_total}$ blocks of \metric{params.block_size_km}~km, coloured by their random split, stratified by sector (seed \metric{params.seed}): training ($n = \metric{blocks.n.train}$), calibration ($n = \metric{blocks.n.calib}$) and test ($n = \metric{blocks.n.test}$). A block belongs to the sector that contains its centre. Map conventions: section~\SImethods, ``Display items''.}
  \label{fig:studyarea}
\end{figure}
}

\newcommand{\SIfigsetup}{%
\begin{figure}[!htbp]
  \centering
  \includegraphics[width=\linewidth]{../figures/out/supp_setup_model.pdf}
  \caption{\textbf{Calibration of the wind-setup model.}
  (a) Coefficient of determination ($R^2$) of the setup model per sector, on calibration blocks, which the fit never used. Labels give $R^2$ and $n$, the numbers of SWOT passes with calibration-block and training-block data in the sector.
  (b) Standard deviation of the de-meaned sector departures of SWOT water-surface elevation from the basin level (grey, no model) and of the residuals after the model (blue), for calibration blocks (filled symbols, upper of each pair) and training blocks (open symbols, lower). Komsomolets Bay has no SWOT open-sea pass and uses the model of the northeastern Kazakh coast. The model regresses the departures on ERA5 wind-stress and pressure features\citep{hersbach2020} (section~\SImethods); in all, \metric{surge.n_passes} SWOT passes and \metric{surge.n_block_passes} block-passes remained after screening. The sector values come from one fit each; no uncertainty is shown.}
  \label{fig:setup}
\end{figure}
}

\newcommand{\SIfigemergence}{%
\begin{figure}[!htbp]
  \centering
  \includegraphics[width=\linewidth]{../figures/out/fig1_emergence.pdf}
  \caption{\textbf{The retreating North Caspian.}
  (a) Year in which each pixel of the emerged band (water in the \metric{params.highstand_label} extent, not in the \metric{params.lowstand_label} extent) became permanently dry, read from the waterline elevation model (\metric{params.dem_res_m}~m; section~\SImethods); grey, no date; light blue, open sea at the low stand. Thin lines and circled numbers, the seven sectors; thick line, the domain. Open circles, ports; triangles, oil and gas platforms and wells; squares, cities (OpenStreetMap)\citep{osm}. The dates inherit the elevation model's errors (figure~\ref{fig:dem}).
  (b) Daily basin level on the DAHITI reference\citep{schwatke2015}, with Hydroweb levels in its gaps (orange)\citep{hydrowebwl,cretaux2011}, and the epochs E1, E2 and E3; the strip under the panel is the colour key of (a).
  (c) Landsat acquisitions per calendar month in the main plan\citep{usgs_l45tm_c2l2,usgs_l7etm_c2l2,usgs_l89oli_c2l2}; the ice-season months are excluded.}
  \label{fig:emergence}
\end{figure}
}

\newcommand{\SIfigmethod}{%
\begin{figure}[!htbp]
  \centering
  \includegraphics[width=\linewidth]{../figures/out/fig2_method.pdf}
  \caption{\textbf{How a falling sea surveys its own seabed.}
  (a) Schematic, not to scale: each waterline, dated by its scene, is a contour of the seabed at the basin level of that day; waterlines are the \metric{params.contour_level} contours of water frequency in \metric{params.bin_width_m}~m level bins (equations~(\ref{eq:si_bins}) and (\ref{eq:si_frequency})).
  (b) Waterlines of one \metric{params.block_size_km}~km training block on the Volga--Ural interfluve, coloured by level (m, DAHITI reference); grey, points removed by the static-contour rule.
  (c) Wind-setup screening on \metric{screening.n_blocks} sampled training blocks: the pre-registered scatter statistic $S(\tau)$ (black; bars, bootstrap standard error) against the threshold $\tau$; horizontal line, one \metric{params.stage2_res_m}~m cell.
  (d) Per sector and level observed on both limbs, the median distance between the E1 and the E2\,+\,E3 shorelines (circles) and between two halves of one epoch (ticks); \metric{h1.e1_vs_e23.n_infinite} of \metric{h1.e1_vs_e23.n_rows} comparisons lacks an E2\,+\,E3 shoreline and is not drawn.}
  \label{fig:method}
\end{figure}
}

\newcommand{\SIfigdem}{%
\begin{figure}[!htbp]
  \centering
  \includegraphics[width=\linewidth]{../figures/out/fig3_dem_accuracy.pdf}
  \caption{\textbf{Seabed topography from waterlines and its accuracy.}
  (a) Waterline elevation model (m, DAHITI reference) in the emerged band, with the \metric{h2.atl06.n_test_blocks} held-out test blocks of \metric{h2.block_size_km}~km (orange).
  (b) Model height against ICESat-2 ATL06 height\citep{atl06_v7} for the $n = \metric{h2.atl06.n_test_points}$ test-block points on bare emerged seabed, in hexagons coloured by point count; reference heights as equation~(\ref{eq:si_lidar_datum}); dashed line, equal heights.
  (c) Left, model minus ICESat-2 against interpolation distance (line, median; shading, interquartile and \metric{params.err_pct_lo}th--\metric{params.err_pct_hi}th percentile ranges); right, NMAD and RMSE per sector of the model (black) and each grid, with test blocks per sector in parentheses.
  (d) Coverage of the \metric{params.conformal_level} conformal intervals on test blocks, overall and per sector, for ATL06 (filled) and ATL08 (open; $n = \metric{h2.atl08.n_test_points}$)\citep{atl08_v7}; bars, test-block bootstrap \metric{params.ci_pct}\% intervals; dashed line, nominal coverage.}
  \label{fig:dem}
\end{figure}
}

\newcommand{\SIfignextmetres}{%
\begin{figure}[!htbp]
  \centering
  \includegraphics[width=\linewidth]{../figures/out/fig5_next_metres.pdf}
  \caption{\textbf{The next metres under the registered correction.}
  (a) Seabed emerged in the domain for declines of \metric{h3.scenarios_text}~m below the 2026 level (\metric{h3.level_2026_m}~m), each surface referenced to its own 2026 shoreline: each grid as used (open)\citep{gebco2026,gebco2022,etopo2022,macferrin2025,tozer2019,amante2009} and the fused surface (filled); bars, \metric{params.conformal_level} split-conformal intervals; grey band, the \metric{h3.n_near_term} near-term declines of the decision rule, whose verdict is printed in each panel.
  (b) The same per sector at \metric{h3.scenario.1_m} and \metric{h3.scenario.2_m}~m: fused estimate (bars) with its interval (lines), and the uncorrected estimate (ticks).
  (c) Exposure in the domain: filled, fused surface; open, grid as used (OpenStreetMap\citep{osm}; ESA WorldCover\citep{worldcover2021}).
  (d) Year in which the primary pathways reach each decline\citep{hoseini2025,samant2023}.
  The fused surface uses the registered error model, which is degenerate in many sectors and did not transfer to the lowest observed levels (section~\SInearterm); it is shown as registered, not as a better estimate.}
  \label{fig:nextmetres}
\end{figure}
}

\newcommand{\SIfigedblocksize}{%
\begin{figure}[!htbp]
  \centering
  \includegraphics[width=\linewidth]{../figures/out/ed_block_size.pdf}
  \caption{\textbf{Sensitivity of the H2 comparison to the validation block size.} For the primary \metric{h2.block_size_km}~km blocks and the \metric{params.block_size_km}~km and \metric{params.block_size_sens_lo_km}~km sensitivities (the number of test blocks under each tick label), upper row ATL06 and lower row ATL08: (a) NMAD and (b) RMSE of the waterline DEM and of each product (matched variant) on the test blocks; (c) (d), the DEM minus product difference in NMAD and RMSE, with the paired test-block bootstrap \metric{params.ci_pct}\% interval (\metric{params.bootstrap_n} resamples); (e), empirical coverage of the \metric{params.conformal_level} conformal intervals with the test-block bootstrap interval, and the nominal level (dashed). H2 population as in figure~\ref{fig:dem}. Each block size has its own random split (same code and seed); the primary verdict rests on the \metric{h2.block_size_km}~km blocks.}
  \label{fig:ed_blocksize}
\end{figure}
}

\newcommand{\SIfigedvegetation}{%
\begin{figure}[!htbp]
  \centering
  \includegraphics[width=\linewidth]{../figures/out/ed_vegetation.pdf}
  \caption{\textbf{The DEM and the products on bare and on vegetated emerged seabed.} Test-block bias (model minus ICESat-2) (a) NMAD (b) and RMSE (c) of each model on the H2 population (filled; emerged seabed with an E3 emergent-vegetation share below \metric{params.veg_cell_share}) and on the vegetated population (open; share of at least \metric{params.veg_cell_share}), for ATL06 (upper marker) and ATL08 (lower marker); products in the matched variant. Point and block counts are printed in the key (vegetated: $n = \metric{h2.veg.atl06.n_test_points}$ ATL06 points in \metric{h2.veg.atl06.n_test_blocks} test blocks). ICESat-2 heights over dense reeds mix canopy and ground, so the vegetated reference is itself uncertain. No uncertainty is drawn; each value pools all test points of a population.}
  \label{fig:ed_vegetation}
\end{figure}
}

\newcommand{\SIfigedscreening}{%
\begin{figure}[!htbp]
  \centering
  \includegraphics[width=\linewidth]{../figures/out/ed_screening.pdf}
  \caption{\textbf{Surge-screening sensitivity.} (a) The pre-registered statistic $S(\tau)$ (black; bootstrap standard error over the \metric{screening.n_blocks} sampled training blocks) and, descriptive only, the mean, 75th and 90th percentiles of the block-bin distances between the shorelines of the two acquisition halves, against the screening threshold $\tau$ on the predicted setup; horizontal line, one \metric{params.stage2_res_m}~m cell. (b) Share of block-bins whose distance exceeds one cell. (c) Number of block-bins with shorelines in both halves. (d) Share of the Landsat acquisitions of the main stage-2 plan whose predicted setup in each sector is at least $\tau$, that is, the acquisitions that screening at $\tau$ would remove there (the number of acquisitions is printed above the panel). The decision rule kept \metric{screening.chosen_tau}. Setup predictions come from the ERA5-driven model\citep{hersbach2020} (figure~\ref{fig:setup}); Komsomolets Bay uses the model of the northeastern Kazakh coast.}
  \label{fig:ed_screening}
\end{figure}
}

\newcommand{\SIfigedhoneshape}{%
\begin{figure}[!htbp]
  \centering
  \includegraphics[width=\linewidth]{../figures/out/ed_h1_blocks_shape.pdf}
  \caption{\textbf{Where each product misplaces the historic shorelines, shape variant.} Median absolute shoreline displacement per \metric{params.block_size_km}~km block, pooled over all (epoch, level) units, for each product shifted by its fitted constant ((a)--(e)). Blocks with fewer than \metric{params.h1_min_block_n} distances are not drawn; the number drawn is printed on each map. Colour key on the figure; light blue, open sea, \metric{params.lowstand_label}; scale bars, \metric{params.fig_scale_bar_km}~km.}
  \label{fig:ed_h1blocks_shape}
\end{figure}
}

\newcommand{\SIfigedhonetotal}{%
\begin{figure}[!htbp]
  \centering
  \includegraphics[width=\linewidth]{../figures/out/ed_h1_blocks_total.pdf}
  \caption{\textbf{Where each product misplaces the historic shorelines, total variant.} As figure~\ref{fig:ed_h1blocks_shape}, for each product on the DAHITI reference through the per-cell geoid correction, without a fitted constant.}
  \label{fig:ed_h1blocks_total}
\end{figure}
}

\newcommand{\SIfigedhexbins}{%
\begin{figure}[!htbp]
  \centering
  \includegraphics[width=\linewidth]{../figures/out/ed_h2_hexbins.pdf}
  \caption{\textbf{Every model against ICESat-2 on the test blocks.} Model height against the ICESat-2 reference height (ellipsoidal height less the EIGEN-6C4 geoid, plus the training-block constant) for every test-block point of the H2 population, ATL06 ((a), $n = \metric{h2.atl06.n_test_points}$) and ATL08 ((b), $n = \metric{h2.atl08.n_test_points}$), in \metric{h2.atl06.n_test_blocks} test blocks of \metric{h2.block_size_km}~km: the waterline DEM and the five products in the matched variant (product plus its training-block offset). Hexagons coloured by point count (logarithmic); dashed line, 1:1. The NMAD and RMSE printed in each panel are those of the test-block accuracy tables.}
  \label{fig:ed_hexbins}
\end{figure}
}
